\section{Модификации генетических операторов}
\label{sec:operators}

\noindent\hspace{1.25cm}\emph{Раздел используется в ЛР~№5--7.}

\subsection{Создание начальной популяции}

Начальная популяция должна быть разнообразной, чтобы ГА с самого начала <<видел>> всю область
поиска. Применяют три способа:
\begin{itemize}
\item \emph{случайная генерация} (<<дробовик>>): каждый ген выбирается равновероятно из допустимых
значений; стандартный способ при отсутствии априорной информации;
\item \emph{фокусировка}: особи генерируются в окрестности известных хороших решений, полученных
эвристикой или предыдущими запусками; ускоряет поиск, но повышает риск преждевременной сходимости;
\item \emph{комбинированный способ}: часть популяции создаётся эвристически, остальные особи~---
случайно.
\end{itemize}
Для перестановочных задач (коммивояжёр) начальные особи~--- случайные перестановки или туры,
построенные жадным методом ближайшего соседа.

\subsection{Селекция}

Оператор селекции формирует из текущей популяции \emph{родительский пул}, используя приспособленность
как единственный источник информации о качестве решений. Методы селекции различаются
\emph{давлением отбора}~--- тем, насколько сильно лучшие особи преобладают в пуле. Слабое давление
замедляет сходимость, сильное ведёт к быстрой потере разнообразия и \emph{преждевременной сходимости}
к локальному экстремуму.

\emph{Пропорциональный отбор} (рулетка, формула (\ref{eq:roulette})) каждой особи сопоставляет сектор
колеса, пропорциональный её приспособленности. Недостатки метода: он требует положительных значений
$F$; зависит от масштаба (прибавление к $F$ большой константы делает выбор почти равновероятным); в
начале поиска одна <<суперособь>> может захватить популяцию, а в конце, когда значения $F$ близки,
давление отбора исчезает. Для устранения этих эффектов применяют \emph{масштабирование}
приспособленности:
\begin{itemize}
\item сдвиг: $F' = F - F_{\min}$ (для максимизации) или $F' = F_{\max} - F$ (для минимизации);
\item линейное: $F' = aF + b$, коэффициенты выбираются так, чтобы среднее не изменилось, а лучшая
особь получала в среднем $C \approx 1,5$--$2$ копии;
\item сигма-отсечение: $F' = \max\bigl(0, F - (\bar F - c\sigma_F)\bigr)$, $c \in [1; 3]$, где
$\sigma_F$~--- стандартное отклонение приспособленности в популяции;
\item степенное: $F' = F^k$, $k > 1$;
\item больцмановское: $F' = \exp(F/T)$ с убывающей <<температурой>> $T$, как при имитации отжига.
\end{itemize}
Реализацию рулетки улучшает \emph{стохастическая универсальная выборка} (SUS, Дж.~Бейкер): колесо
вращают один раз, а $N$ указателей расставлены с шагом $1/N$. Число копий каждой особи тогда
отличается от ожидаемого $Np_i$ не более чем на единицу.

\emph{Ранжирование.} Особи сортируются по приспособленности, и вероятность выбора зависит только от
ранга $i$ ($i = 1$~--- лучшая), а не от самого значения $F$. При линейном ранжировании
\begin{equation}
\label{eq:rank}
p_i = \frac{1}{N}\left(a - (a - b)\frac{i - 1}{N - 1}\right), \qquad 1 \le a \le 2, \quad b = 2 - a,
\end{equation}
где $a$~--- ожидаемое число копий лучшей особи. Метод не чувствителен к масштабу и знаку $F$ и сохраняет
давление отбора, когда значения приспособленности близки.

\emph{Отбор усечением} (truncation): в пул допускаются только $T\cdot 100\,\%$ лучших особей
($10\,\% \le T \le 50\,\%$), среди которых родители выбираются равновероятно. Применяется в основном
для больших популяций.

\emph{Турнирный отбор}: из популяции случайно выбираются $k$ особей, и лучшая из них становится
родителем; процедура повторяется $N$ раз. Размер турнира $k$ ($2 \le k \le N$) прямо регулирует
давление отбора: вероятность выбора особи ранга $i$ равна
$\bigl((N - i + 1)^k - (N - i)^k\bigr)/N^k$. Турнир не зависит от масштаба $F$, легко распараллеливается
и потому является самым распространённым методом в современных ГА. Сравнение методов приведено на
рисунке~\ref{fig:selection}.

\begin{figure}[!htb]
\centering
\includegraphics[width=\textwidth]{\figdir/selection.png}
\caption{Вероятности выбора особей при разных методах селекции (популяция из 10 особей)}
\label{fig:selection}
\end{figure}

\emph{Локальный отбор} выполняется среди соседей особи в пространственно структурированной популяции
(линейной, двумерной с 4- или 8-связностью). Малые окрестности замедляют распространение хороших
решений и помогают сохранить разнообразие.

После формирования пула выбираются \emph{пары для скрещивания}:
\begin{itemize}
\item \emph{панмиксия}~--- оба родителя выбираются случайно; универсальный метод;
\item \emph{селективный выбор}~--- родителями могут стать только особи с приспособленностью не ниже
средней; ускоряет сходимость, но опасен для многоэкстремальных задач;
\item \emph{инбридинг}~--- второй родитель выбирается как ближайший к первому (генотипически, по
расстоянию Хэмминга, или фенотипически, по расстоянию в пространстве решений); способствует
локальному поиску внутри <<ниши>>;
\item \emph{аутбридинг}~--- пары образуют наиболее далёкие особи; поддерживает разнообразие.
\end{itemize}
Комбинация инбридинга и аутбридинга применяется при поиске нескольких экстремумов.

\subsection{Скрещивание двоичных хромосом}

Оператор скрещивания (кроссинговер, рекомбинация) создаёт потомков, комбинируя гены двух родителей.
Он применяется к паре с вероятностью $P_c$; в противном случае потомки~--- копии родителей. Основные
двоичные операторы показаны на рисунке~\ref{fig:crossover}:
\begin{itemize}
\item \emph{одноточечный}: случайно выбирается точка $k \in \{1, \ldots, L - 1\}$, и родители
обмениваются частями хромосом после неё;
\item \emph{многоточечный}: выбираются $m$ точек, и родители обмениваются фрагментами через один;
двухточечный кроссинговер меньше разрушает схемы, фиксированные позиции которых находятся на концах
хромосомы;
\item \emph{однородный}: генерируется случайная двоичная маска, и каждый ген первого потомка берётся из
первого родителя там, где в маске 1, и из второго там, где 0 (второй потомок~--- наоборот). Каждая
позиция является потенциальной точкой обмена, поэтому однородный кроссинговер сильнее всего
перемешивает гены и не зависит от их взаимного расположения.
\end{itemize}

\begin{figure}[!htb]
\centering
\begin{tikzpicture}[>=Stealth]
\def\A{cblue!22}\def\B{corange!30}
% одноточечный
\node[font=\small\bfseries] at (2.6, 1.3) {одноточечный, $k = 4$};
\bitrow[0.4]{0}{$A$}{0/\A,1/\A,1/\A,0/\A,1/\A,0/\A,0/\A,1/\A,1/\A,0/\A}
\bitrow[0.4]{-0.55}{$B$}{1/\B,0/\B,0/\B,1/\B,1/\B,1/\B,0/\B,0/\B,1/\B,1/\B}
\bitrow[0.4]{-1.45}{$A'$}{0/\A,1/\A,1/\A,0/\A,1/\B,1/\B,0/\B,0/\B,1/\B,1/\B}
\bitrow[0.4]{-2.0}{$B'$}{1/\B,0/\B,0/\B,1/\B,1/\A,0/\A,0/\A,1/\A,1/\A,0/\A}
\cutline[0.4]{4}{0}{-0.55}
\draw[->, cgray] (2.6, -0.85) -- (2.6, -1.15);
% двухточечный
\node[font=\small\bfseries] at (8.2, 1.3) {двухточечный, $k = 3, 7$};
\bitrow[6.0]{0}{$A$}{0/\A,1/\A,1/\A,0/\A,1/\A,0/\A,0/\A,1/\A,1/\A,0/\A}
\bitrow[6.0]{-0.55}{$B$}{1/\B,0/\B,0/\B,1/\B,1/\B,1/\B,0/\B,0/\B,1/\B,1/\B}
\bitrow[6.0]{-1.45}{$A'$}{0/\A,1/\A,1/\A,1/\B,1/\B,1/\B,0/\B,1/\A,1/\A,0/\A}
\bitrow[6.0]{-2.0}{$B'$}{1/\B,0/\B,0/\B,0/\A,1/\A,0/\A,0/\A,0/\B,1/\B,1/\B}
\cutline[6.0]{3}{0}{-0.55}\cutline[6.0]{7}{0}{-0.55}
\draw[->, cgray] (8.2, -0.85) -- (8.2, -1.15);
% однородный
\node[font=\small\bfseries] at (13.8, 1.3) {однородный};
\bitrow[11.6]{0.62}{$M$}{1/white,0/white,0/white,1/white,1/white,0/white,1/white,0/white,0/white,1/white}
\bitrow[11.6]{0}{$A$}{0/\A,1/\A,1/\A,0/\A,1/\A,0/\A,0/\A,1/\A,1/\A,0/\A}
\bitrow[11.6]{-0.55}{$B$}{1/\B,0/\B,0/\B,1/\B,1/\B,1/\B,0/\B,0/\B,1/\B,1/\B}
\bitrow[11.6]{-1.45}{$A'$}{0/\A,0/\B,0/\B,0/\A,1/\A,1/\B,0/\A,0/\B,1/\B,0/\A}
\bitrow[11.6]{-2.0}{$B'$}{1/\B,1/\A,1/\A,1/\B,1/\B,0/\A,0/\B,1/\A,1/\A,1/\B}
\draw[->, cgray] (13.8, -0.85) -- (13.8, -1.15);
\end{tikzpicture}
\caption{Двоичные операторы скрещивания: $A$, $B$~--- родители, $A'$, $B'$~--- потомки,
$M$~--- маска однородного кроссинговера}
\label{fig:crossover}
\end{figure}

\subsection{Мутация двоичных хромосом}

Мутация вносит в популяцию новые значения генов, которых могло не быть ни у одного родителя, и
предотвращает вырождение популяции. Для двоичного кодирования применяются
(рисунок~\ref{fig:mutation}):
\begin{itemize}
\item \emph{побитовая мутация}~--- каждый бит инвертируется с вероятностью $P_m$ (типично $1/L$, т.\,е. в
среднем один бит на хромосому);
\item \emph{инверсия}~--- выбираются две позиции, и отрезок хромосомы между ними записывается в
обратном порядке; порядок генов меняется, а их состав (число нулей и единиц) сохраняется;
\item \emph{обмен (перестановка)}~--- два случайных гена меняются местами.
\end{itemize}
Инверсия и обмен не создают новых аллелей: если во всей популяции в некоторой позиции стоит 0, эти
операторы могут перенести туда единицу из другой позиции хромосомы, но не гарантируют этого. Поэтому
для числовых задач с двоичным кодированием их используют вместе с побитовой мутацией.

\begin{figure}[!htb]
\centering
\begin{tikzpicture}
\def\A{cblue!22}\def\M{cgreen!45}
\node[font=\small\bfseries] at (2.5, 0.75) {побитовая, позиция 4};
\bitrow[0.5]{0}{до}{0/\A,1/\A,1/\A,0/\M,1/\A,0/\A,0/\A,1/\A,1/\A,0/\A}
\bitrow[0.5]{-0.6}{после}{0/\A,1/\A,1/\A,1/\M,1/\A,0/\A,0/\A,1/\A,1/\A,0/\A}
\node[font=\small\bfseries] at (7.8, 0.75) {инверсия, позиции 3--7};
\bitrow[5.8]{0}{}{0/\A,1/\A,1/\M,0/\M,1/\M,0/\M,0/\M,1/\A,1/\A,0/\A}
\bitrow[5.8]{-0.6}{}{0/\A,1/\A,0/\M,0/\M,1/\M,0/\M,1/\M,1/\A,1/\A,0/\A}
\node[font=\small\bfseries] at (13.1, 0.75) {обмен, позиции 1 и 5};
\bitrow[11.1]{0}{}{0/\M,1/\A,1/\A,0/\A,1/\M,0/\A,0/\A,1/\A,1/\A,0/\A}
\bitrow[11.1]{-0.6}{}{1/\M,1/\A,1/\A,0/\A,0/\M,0/\A,0/\A,1/\A,1/\A,0/\A}
\end{tikzpicture}
\caption{Операторы мутации двоичной хромосомы (изменённые гены выделены)}
\label{fig:mutation}
\end{figure}

\subsection{Вещественное кодирование}
\label{sec:real}

При \emph{вещественном кодировании} особь~--- вектор $\mathbf{x} = (x_1, \ldots, x_n)$ самих
параметров задачи. Оно не ограничено сеткой двоичного кода, не требует декодирования, естественно для
непрерывных задач и позволяет использовать операторы, учитывающие метрику пространства. Операторы
скрещивания для вещественных векторов~\cite{skobtsov2018, michalewicz1996}:
\begin{itemize}
\item \emph{дискретная рекомбинация}~--- каждая компонента потомка берётся из первого или второго
родителя с равной вероятностью (аналог однородного кроссинговера);
\item \emph{промежуточная рекомбинация}~--- компоненты потомка лежат между значениями родителей или
в их окрестности:
\begin{equation}
\label{eq:intermediate}
c_i = p_i^1 + \alpha_i\bigl(p_i^2 - p_i^1\bigr), \qquad \alpha_i \in [-d; 1 + d],
\end{equation}
где $\alpha_i$ выбирается случайно для каждой компоненты; при $d = 0$ потомок лежит внутри
гиперпрямоугольника родителей, при $d > 0$ (обычно $d = 0,25$)~--- и вне его;
\item \emph{линейная рекомбинация}~--- то же, но $\alpha$ одно для всех компонент: потомок лежит на
прямой, проходящей через родителей; частный случай $\alpha \in [0; 1]$~--- \emph{арифметический}
кроссинговер $\mathbf{c} = \lambda\mathbf{p}^1 + (1 - \lambda)\mathbf{p}^2$;
\item \emph{BLX-$\alpha$} (blend crossover, Л.~Эшелман и Д.~Шаффер~\cite{eshelman1993})~---
компонента потомка выбирается равномерно из интервала, расширенного на долю $\alpha$ расстояния между
родителями:
\begin{equation}
\label{eq:blx}
c_i \sim U\bigl[\min(p_i^1, p_i^2) - \alpha d_i;\ \max(p_i^1, p_i^2) + \alpha d_i\bigr], \qquad
d_i = |p_i^1 - p_i^2| .
\end{equation}
BLX-$\alpha$ совпадает с промежуточной рекомбинацией при $d = \alpha$. Значение $\alpha \approx 0,5$
сохраняет разброс популяции в среднем неизменным, $\alpha = 0$ сжимает популяцию, большие $\alpha$
приближают поиск к случайному (рисунок~\ref{fig:real-crossover}).
\end{itemize}

\begin{figure}[!htb]
\centering
\includegraphics[width=0.62\textwidth]{\figdir/real_crossover.png}
\caption{Области возможных потомков при вещественном скрещивании двух родителей}
\label{fig:real-crossover}
\end{figure}

Мутация вещественной особи~--- прибавление к гену случайного \emph{шага мутации}:
\begin{itemize}
\item \emph{равномерная}: ген заменяется случайным значением из допустимого отрезка $[a_i; b_i]$;
\item \emph{гауссова}: $x_i' = x_i + \mathcal{N}(0, \sigma^2)$; малый шаг даёт точность, большой~---
способность выходить из локальных экстремумов. Поэтому шаг уменьшают по ходу поиска, например
линейно: $\sigma(t) = \sigma_0(b - a)(1 - t/T)$, где $T$~--- число поколений;
\item \emph{неравномерная мутация} З.~Михалевича~\cite{michalewicz1996}: $x_i' = x_i + \Delta(t, b_i - x_i)$
или $x_i' = x_i - \Delta(t, x_i - a_i)$ с равной вероятностью, где
\begin{equation}
\label{eq:nonuniform}
\Delta(t, y) = y\left(1 - r^{(1 - t/T)^{\beta}}\right),
\end{equation}
$r \sim U[0; 1]$, $\beta \approx 2$. В начале поиска шаг может достигать границ области, к концу он
стремится к нулю, и мутация выполняет локальную доводку.
\end{itemize}
Если после мутации или BLX-$\alpha$ ген выходит за пределы допустимого отрезка, его возвращают на
границу (отсечение) или отражают внутрь области.

\subsection{Сокращение промежуточной популяции}

После скрещивания и мутации получается промежуточная популяция из родителей и потомков, численность
которой больше $N$. Стратегия сокращения (редукции) определяет, кто перейдёт в следующее поколение:
\begin{itemize}
\item \emph{чистая замена} (поколенческая схема): $N$ потомков полностью заменяют родителей; каждая
особь живёт одно поколение, и лучшее решение может быть потеряно;
\item \emph{элитарная схема} ($\mu + \lambda$): из объединения родителей и потомков отбираются $N$
лучших; найденное лучшее решение никогда не теряется, но разнообразие убывает быстрее. Частный случай~---
\emph{элитизм}: в новое поколение без изменений переносятся 1--2 лучшие особи, остальные места
занимают потомки;
\item \emph{равномерная случайная замена}: $N$ особей выбираются из объединения равновероятно;
\item \emph{пропорциональная редукция}: особи выбираются из объединения рулеткой по приспособленности;
\item \emph{отбор вытеснением} (crowding): потомок заменяет наиболее похожую на него особь, что
сохраняет разнообразие и позволяет поддерживать несколько экстремумов;
\item \emph{схема Больцмана}: потомок заменяет родителя с вероятностью, зависящей от разности их
приспособленностей и убывающей температуры;
\item \emph{локальная замена}: в пространственно структурированной популяции потомок заменяет особь из
окрестности своего родителя.
\end{itemize}
В \emph{стационарном} ГА (steady-state) за итерацию создаётся один--два потомка, которые заменяют
худшие особи; поколения в явном виде не выделяются.

\subsection{Настройка параметров и методика эксперимента}
\label{sec:ga-method}

Работа ГА определяется балансом двух тенденций: \emph{исследования} (exploration)~--- поиска в новых
областях, которое обеспечивают мутация, большая популяция и слабое давление отбора, и
\emph{использования} (exploitation)~--- уточнения найденных хороших решений, которое обеспечивают
селекция, элитизм и скрещивание близких особей. Перекос в сторону исследования превращает ГА в
случайный поиск, перекос в сторону использования ведёт к преждевременной сходимости.
Ориентировочные значения параметров приведены в таблице~\ref{tab:ga-params}.

\begin{table}[!htb]
\caption{Типичные значения параметров генетического алгоритма}
\label{tab:ga-params}
\centering\small
\begin{tabularx}{\textwidth}{>{\raggedright}p{4.6cm}>{\raggedright}p{3.1cm}X}
\toprule
Параметр & Диапазон & Влияние \\
\midrule
Размер популяции $N$ & 20--200 & главный фактор надёжности; растёт с числом переменных и числом экстремумов \\
Вероятность скрещивания $P_c$ & 0,6--0,95 & ускоряет сходимость; при $P_c = 0$ ГА сводится к мутационному поиску \\
Вероятность мутации $P_m$ & $1/L$; 0,001--0,01 на бит & без мутации популяция вырождается, при большой~--- случайный поиск \\
Размер турнира $k$ & 2--5 & давление отбора \\
Число элитных особей & 1--2 & монотонность лучшего решения \\
Длина хромосомы $L$ & формула (\ref{eq:chrom-length}) & предельная точность и размер пространства поиска \\
\bottomrule
\end{tabularx}
\end{table}

Результат одного запуска ГА случаен, поэтому выводы делают по серии из 20--50 независимых запусков с
разными начальными значениями генератора случайных чисел. Для каждой настройки фиксируют:
\begin{itemize}
\item долю \emph{успешных} запусков~--- найдено решение с заданной точностью относительно известного
оптимума;
\item медиану и худшее значение найденного экстремума;
\item \emph{скорость}~--- медиану числа поколений или вычислений функции до первого успеха.
\end{itemize}
При исследовании параметра меняют только его, оставляя остальные неизменными. Сравнение разных
алгоритмов корректно только при одинаковом бюджете вычислений функции.

\subsection{Векторная реализация операторов на NumPy}
\label{sec:ga-numpy}

Популяцию из $N$ двоичных хромосом длины $L$ удобно хранить массивом \code{pop} формы $(N, L)$ типа
\code{int8}, приспособленности~--- вектором \code{F}. Тогда каждый оператор~--- одна-две строки без
циклов по особям:
\begin{lstlisting}
idx = rng.choice(N, N, p=F / F.sum())                 # рулетка (F > 0)
cand = rng.integers(0, N, (N, k))                     # турнир: k случайных особей,
idx = cand[np.arange(N), F[cand].argmax(axis=1)]      # побеждает лучшая
a, b = pop[idx][0::2], pop[idx][1::2]                 # пары родителей
cut = rng.integers(1, L, (len(a), 1))                 # одноточечный кроссинговер
mask = (np.arange(L) >= cut) & (rng.random((len(a), 1)) < pc)
kids = np.vstack([np.where(mask, b, a), np.where(mask, a, b)])
kids ^= (rng.random(kids.shape) < pm).astype(np.int8) # побитовая мутация
t = lo + kids @ 2 ** np.arange(L - 1, -1, -1) * (hi - lo) / (2 ** L - 1)   # декодирование (5.2)
\end{lstlisting}
Для однородного кроссинговера маска генерируется как \code{rng.random(a.shape) < 0.5}, для BLX-$\alpha$
потомок~--- \code{rng.uniform(np.minimum(a, b) - alpha * d, np.maximum(a, b) + alpha * d)}.

Готовые библиотеки используются для проверки своей реализации, а не вместо неё. В DEAP~\cite{fortin2012}
операторы собраны в модуле \code{deap.tools} (\code{selTournament}, \code{cxTwoPoint}, \code{cxUniform},
\code{mutFlipBit}, \code{cxBlend}, \code{mutGaussian}), в PyGAD задаются параметрами
\code{parent\_selection\_type}, \code{crossover\_type}, \code{mutation\_type}. Эталонный экстремум для
оценки точности находят детерминированными методами: \code{scipy.optimize.minimize\_scalar} на
мелкой сетке для функции одной переменной и \code{scipy.optimize.differential\_evolution} для функций
двух переменных.
